% !TeX root = ../../Book2.tex
\section{热半群}
\label{b2:sec:1204}

\subsection{半群性质}
\label{b2:subsec:120401}

将时刻 \(t\) 的解视为对初值的作用，记
\[
 S(t)f\coloneqq H_{\kappa t}*f\quad(t>0),
 \quad S(0)f=f.
\]
这一族算子称为\term[re-banqun]{热半群}[heat semigroup]。先演化 \(s\) 再演化 \(t\)，应与直接演化 \(s+t\) 相同；核公式使这一预期成为可以验证的等式。
\symidx{\(S(t)\)：全空间热半群}

\begin{proposition}\label{b2:prop:12-semigroup}
对 \(s,t\geq0\)，有 \(S(t)S(s)=S(t+s)\)。当 \(s,t>0\) 时，等价的核关系为 \(H_{\kappa t}*H_{\kappa s}=H_{\kappa(t+s)}\)。这些等式在全部 \(L^p\)、\(1\leq p\leq\infty\)，以及有界一致连续函数空间上成立。
\end{proposition}
\begin{proof}
两个热核均可积，卷积定理给出变换
\[
 e^{-4\pi^2\kappa t|\xi|^2}e^{-4\pi^2\kappa s|\xi|^2}
 =e^{-4\pi^2\kappa(t+s)|\xi|^2}.
\]
Fourier 唯一性使核相等。对一般 \(L^p\) 初值，卷积结合律可先对有界紧支集函数由 Fubini 证明，再由 Young 不等式和稠密性推广；\(p=\infty\) 时，双重积分的绝对值直接由 \(\|f\|_\infty\) 乘两个核的质量控制，Fubini 仍适用。零参数时是恒等算子。
\end{proof}

区间上的对应算子记为 \(S_I(t)\)，由式\eqref{b2:eq:12-interval-expansion}按每个特征系数相乘，也有同一半群关系。全空间借助平移不变的核，区间借助边界适配的特征基；两者组织时间的方式相同。

\subsection{压缩性}
\label{b2:subsec:120402}

\begin{proposition}[范数收缩与强连续性]\label{b2:prop:12-contraction}
全空间热半群满足 \(\|S(t)f\|_p\leq\|f\|_p\)，\(1\leq p\leq\infty\)。对 \(p<\infty\)，映射 \(t\mapsto S(t)f\) 在 \([0,\infty)\) 上按 \(L^p\) 范数连续。这个连续性也在 \(C_0(\mathbb R^d)\) 与有界一致连续函数空间上成立，但不在整个 \(L^\infty\) 上成立。
\end{proposition}
\begin{proof}
热核质量为一，Young 不等式给出范数界。零时连续性已由逼近恒等证明。若 \(h>0\)，半群关系给出
\[
 \|S(t+h)f-S(t)f\|_p\leq\|S(h)f-f\|_p.
\]
对左侧时间差，用 \(S(t)f-S(t-h)f=S(t-h)(S(h)f-f)\) 得同样控制。于是零时连续性延伸到全部时间。

对 \(C_0\) 还须检查不变性。连续且在无穷远消失的函数可由连续紧支集函数一致逼近。后者与热核卷积在无穷远消失：初值支集有界，核在 \(|x|\to\infty\) 时于该支集上一致趋零。范数界让这项性质传到一般 \(C_0\) 函数；连续性和一致连续性由核的平移连续性得到。有界一致连续空间的不变性则直接由 \(|S(t)f(x+h)-S(t)f(x)|\leq\sup_y|f(y+h)-f(y)|\) 得出。

最后取 \(f=\mathbf1_{\{x_1>0\}}\)。一维积分公式表明对每个 \(t>0\)，\(\|S(t)f-f\|_\infty=1/2\)：在界面两侧趋近界面，误差趋于 \(1/2\)，而远离界面时严格更小。因此全体 \(L^\infty\) 的强连续性失败。
\end{proof}

范数不增不等于每个非零函数都按一个共同的指数衰减。在 \(L^2(\mathbb R^d)\) 上，乘子 \(e^{-4\pi^2\kappa t|\xi|^2}\) 的本性上确界为一，因此 \(\|S(t)\|_{2\to2}=1\)。任意低频附近的数据衰减都可以很慢。另一方面，对每个固定 \(f\in L^2\)，这个乘子在几乎每个频率处趋零，控制收敛给出 \(\|S(t)f\|_2\to0\)。逐个初值衰减与统一算子范数衰减不能混为一谈。

\ProseParagraphBreak

Dirichlet 与 Neumann 区间热算子也在全部 \(L^p\) 上收缩。Neumann 核非负且对两变量积分都为一；Dirichlet 核非负且两类积分均不超过一。后一个断言可由初值一的解与常数一比较得到，再用核对称性换变量。对 \(1<p<\infty\)，补上一份零函数的质量后应用 Jensen，不等式为 \(|S_I(t)f|^p\leq S_I(t)|f|^p\)，再积分即得；两个端点直接估计。Dirichlet 核的非负性由上一节正性保持与连续核得到，若某个核值为负，以支集集中于该点附近的非负初值测试就会矛盾。

\subsection{生成元观点}
\label{b2:subsec:120403}

半群给出了有限时间的变化，生成元记录零时的一阶变化。设 \((T(t))_{t\geq0}\) 是 Banach 空间 \(X\) 上满足 \(T(0)=I\)、\(T(t+s)=T(t)T(s)\) 的有界线性算子族，并且 \(T(t)x\to x\) 当 \(t\downarrow0\) 对每个 \(x\in X\) 成立。这样的族称为\term[qiang-lianxu-banqun]{强连续半群}[strongly continuous semigroup]。定义
\begin{equation}\label{b2:eq:12-generator-def}
 D(L)\coloneqq\left\{\,f\in X\;\middle|\;
 \lim_{t\downarrow0}\dfrac{T(t)f-f}{t}\text{ 在 }X\text{ 中存在}\,\right\},
 \quad Lf\coloneqq\lim_{t\downarrow0}\dfrac{T(t)f-f}{t},
\end{equation}
称 \(L\) 为半群的\term[shengchengyuan]{生成元}[generator]。这个极限必须在所选空间的范数中成立，不能用逐点导数替代定义域条件。

\begin{theorem}[热半群的平方可积生成元]\label{b2:thm:12-generator-l2}
在 \(X=L^2(\mathbb R^d)\) 上，热半群的生成元为
\[
 L=\kappa\Delta,
 \quad D(L)=H^2(\mathbb R^d).
\]
对 \(\lambda>0\)，其逆算子有
\[
 (\lambda I-L)^{-1}f=\int_0^\infty e^{-\lambda t}S(t)f\,\mathrm{d}t,
 \quad\|(\lambda I-L)^{-1}\|_{2\to2}\leq\lambda^{-1}.
\]
\end{theorem}
\begin{proof}
仍令 \(a(\xi)=4\pi^2\kappa|\xi|^2\)。若 \(a\widehat f\in L^2\)，则
\[
 \left|\dfrac{e^{-ta}-1}{t}\right|\leq a,
 \quad\dfrac{e^{-ta}-1}{t}\widehat f\longrightarrow-a\widehat f
 \quad\text{于 }L^2.
\]
控制收敛与 Plancherel 说明 \(f\in D(L)\)，且 \(Lf=\kappa\Delta f\)。反过来，若差商在 \(L^2\) 中收敛，变换也收敛；取一个几乎处处收敛的子列，与乘子逐点极限比较，得到极限必为 \(-a\widehat f\)，故 \(a\widehat f\in L^2\)。由\cref{b2:thm:04-weak-derivative}，条件 \(f,|\xi|^2\widehat f\in L^2\) 等价于 \(f\in H^2\)：各二次单项式均受 \(|\xi|^2\) 控制，一次单项式则由 \(1+|\xi|^2\) 控制。

最后，指数权重使算子积分在 \(L^2\) 范数中收敛。变换后可逐频率积分，乘子为 \(\int_0^\infty e^{-t(\lambda+a)}\,\mathrm{d}t=(\lambda+a)^{-1}\)。其模不超过 \(\lambda^{-1}\)，而 \(a/(\lambda+a)\leq1\)，故积分的像属于生成元定义域。乘上 \(\lambda+a\) 即验证双侧逆。
\end{proof}

区间的生成元为 \(-A\)，定义域仍是\cref{b2:def:10-sl-operator}的边界定义域；系数差商与同一控制论证证明这一点。负号由 \(u_t+Au=0\) 决定。相同的微分表达式 \(\kappa\partial_x^2\)，配上不同边界定义域，就产生不同的半群和不同的长期平衡。

\subsection{\texorpdfstring{\(C_0\)}{C0} 半群的预告}
\label{b2:subsec:120404}

强连续半群也称 \(C_0\) 半群；这里的下标表示零时强连续，与函数空间 \(C_0(\mathbb R^d)\) 表示无穷远消失不是同一个定义。若还满足 \(\|T(t)\|\leq1\)，称为\term[yasuo-banqun]{压缩半群}[contraction semigroup]；此处允许范数等于一，不要求压缩映射不动点定理中的严格压缩常数。

\ProseParagraphBreak

在热流中，半群先由显式核构造，再识别生成元。更一般的扩散模型常没有可写出的核，问题顺序会反过来：先给闭算子及其定义域，再由解算子的估计构造半群。\appref{b2:app:D}将介绍这条抽象路线。当前结果只需读者掌握三点：时间连续性依赖所选空间，生成元必须连同定义域一起给出，指数演化的记号并不自动允许反向时间。

\ProseParagraphBreak

反向时间会放大高频。在全空间 \(L^2\) 中，若想由时刻 \(T>0\) 的值恢复初值，频率公式要求乘以 \(e^{4\pi^2\kappa T|\xi|^2}\)。该乘子无界，故即使正时间演化是单射，逆问题也不具有同一范数下的连续依赖。温度演化的稳定性与反演初温的稳定性是两个不同的问题。
